\documentclass[10pt,letterpaper]{article}
\usepackage{cornell-notes}

\title{Chapter 6: The Transport Layer}
\author{}
\date{}

\setDocTitle{Chapter 6: The Transport Layer}
\setDocSubtitle{Computer Networks}
\setDocVersion{v1.0}
\setDocStatus{Study Companion}
\setDocOwner{Computer Networks Cornell Notes}
\setDocAuthor{}
\setDocDate{\today}
\setDocSummary{Cornell-format study and review notes for Chapter 6: The Transport Layer.}

\setCornellCollection{Computer Networks Cornell Notes}
\setCornellUnitType{Chapter}
\setCornellUnitNumber{6}
\setCornellUnitTitle{The Transport Layer}
\setCornellCompanionLabel{Study and review companion}

\hypersetup{pdftitle={Chapter 6: The Transport Layer},pdfsubject={Cornell notes},pdfauthor={}}

\begin{document}
\maketitle
\notesection{Chapter overview}
\begin{overviewbox}
\textbf{Learning objectives}\begin{itemize}
  \item Explain transport service primitives, sockets, and process addressing.
  \item Reason about connection establishment, release, flow/error control, multiplexing, and crash recovery.
  \item Compare UDP, RPC, RTP, and TCP service choices.
  \item Trace TCP headers, state transitions, windows, timers, and congestion behavior.
  \item Relate network performance to host design, bandwidth-delay product, and delay-tolerant operation.
\end{itemize}
\textbf{Prerequisite concepts}\quad Network-layer datagrams, sliding windows, congestion, process abstractions, and basic probability.\par
\textbf{Key themes}\quad End-to-end service; ports; connection state; reliability; receiver flow control; congestion response; performance; disruption tolerance.\par
\textbf{Named mechanisms and standards}\quad Berkeley sockets, three-way handshake, graceful release, UDP, RPC, RTP/RTCP, TCP sequence and acknowledgment fields, retransmission timers, sliding windows, congestion windows, header compression, and the Bundle Protocol.
\end{overviewbox}

\notesection{Cornell cue-and-note record}

\begin{longtable}{@{}L{0.245\textwidth}|L{0.695\textwidth}@{}}
\rowcolor{CNavy}\color{white}\bfseries Cue / recall prompt &
\color{white}\bfseries Notes \\
\endfirsthead
\rowcolor{CNavy}\color{white}\bfseries Cue / recall prompt &
\color{white}\bfseries Notes (continued) \\
\endhead
\hline
\endfoot
\cellcolor{CCue}\textbf{6.1 Transport service} & \begin{minipage}[t]{\linewidth}\raggedright \begin{itemize}\item \textbf{6.1.1 Services to upper layers} Transport isolates applications from network details and provides process-to-process delivery. It can add connection management, reliability, ordering, and flow regulation above a less reliable network.\item \textbf{6.1.2 Service primitives} A connection-oriented interface commonly supports listening, connecting, accepting, sending, receiving, and disconnecting. Blocking behavior and buffering affect how applications observe the service.\item \textbf{6.1.3 Berkeley sockets} Sockets expose endpoints through calls such as socket, bind, listen, accept, connect, send, receive, and close. The interface is local; TCP or UDP messages are the peer protocols.\item \textbf{6.1.4 File-server example} A server binds a known endpoint, listens, accepts connections, reads requests, returns data, and closes. A client creates a socket, connects, exchanges data, and terminates.\end{itemize}\end{minipage} \\ \hline
\cellcolor{CCue}\textbf{6.2 Elements of transport protocols} & \begin{minipage}[t]{\linewidth}\raggedright \begin{itemize}\item \textbf{6.2.1 Addressing} Transport service access points identify processes, commonly with ports. Servers use stable or discoverable endpoints; clients often receive temporary endpoints.\item \textbf{6.2.2 Connection establishment} Delayed duplicate packets can mimic a new request. Sequence-number protection and a three-way exchange let both sides confirm fresh state and negotiate starting values.\item \textbf{6.2.3 Connection release} Asymmetric release may lose in-flight data; symmetric release treats each direction separately. The two-army problem shows that no finite exchange over an unreliable channel creates perfect common knowledge.\item \textbf{6.2.4 Error and flow control} Checksums, sequence numbers, acknowledgments, timers, buffers, and variable windows provide end-to-end correctness and prevent receiver overrun.\item \textbf{6.2.5 Multiplexing} Upward multiplexing lets many transport connections share one network endpoint; downward multiplexing can distribute one transport demand across multiple network connections.\item \textbf{6.2.6 Crash recovery} A crash may erase connection state while old packets persist. Recovery requires stable information or application-level reconciliation; transport alone cannot always decide whether an operation took effect.\end{itemize}\end{minipage} \\ \hline
\cellcolor{CCue}\textbf{6.3 Congestion control} & \begin{minipage}[t]{\linewidth}\raggedright Transport congestion control limits pressure on the network; receive-window flow control limits pressure on the destination.\par\begin{itemize}\item \textbf{6.3.1 Desirable allocation} An allocation should be efficient and fair. Additive increase and multiplicative decrease move competing flows toward a useful operating point under feedback.\item \textbf{6.3.2 Regulating send rate} Window-based control limits outstanding data; rate-based control paces sending. Feedback may be inferred from loss or delay or carried explicitly by the network.\item \textbf{6.3.3 Wireless issues} Wireless corruption and route changes can cause loss unrelated to congestion. Treating every loss as congestion may reduce rate unnecessarily, yet hiding all loss can also defeat congestion signals.\end{itemize}\end{minipage} \\ \hline
\cellcolor{CCue}\textbf{6.4 UDP} & \begin{minipage}[t]{\linewidth}\raggedright \begin{itemize}\item \textbf{6.4.1 Introduction to UDP} UDP adds source and destination ports, length, and checksum to IP. It preserves message boundaries but does not establish a connection, retransmit, order, or regulate congestion by itself.\item \textbf{6.4.2 Remote procedure call} RPC represents a call as request and reply messages. Stubs marshal parameters, but failures make execution semantics difficult: a timeout cannot reveal whether a remote operation never ran, ran once, or ran before a reply was lost.\item \textbf{6.4.3 Real-time transport} RTP carries timestamps, sequence numbers, payload type, and source identity for media; RTCP reports reception and synchronization information. RTP supports application timing but does not reserve resources or guarantee delivery.\end{itemize}\end{minipage} \\ \hline
\cellcolor{CCue}\textbf{6.5 TCP} & \begin{minipage}[t]{\linewidth}\raggedright \begin{itemize}\item \textbf{6.5.1 Introduction} TCP provides a reliable, ordered, full-duplex byte stream over an internetwork that may lose, duplicate, reorder, or delay packets.\item \textbf{6.5.2 Service model} Applications use sockets and connections identified by endpoint pairs. TCP is a byte stream, so application message boundaries are not preserved.\item \textbf{6.5.3 Protocol} Bytes receive sequence numbers; segments carry byte ranges; cumulative acknowledgments report the next expected byte. Retransmission and reassembly hide many network faults.\item \textbf{6.5.4 Segment header} Ports identify processes; sequence and acknowledgment numbers track bytes; flags control state; window advertises receive capacity; checksum protects header and data; options negotiate extensions.\item \textbf{6.5.5 Establishment} The three-way handshake exchanges initial sequence numbers and confirms that both directions are ready, reducing confusion from delayed duplicates.\item \textbf{6.5.6 Release} Each direction closes independently using control flags and acknowledgments. Time-wait state prevents old segments from being mistaken for a later connection.\item \textbf{6.5.7 Connection modeling} A finite-state machine captures transitions among listen, synchronization, established, closing, and time-wait states in response to calls, segments, and timeouts.\item \textbf{6.5.8 Sliding window} The effective send limit reflects the receiver's advertised window and the congestion limit. Cumulative acknowledgments advance the window; buffering handles retransmission and in-order delivery.\item \textbf{6.5.9 Timer management} Retransmission timeout adapts to measured round-trip time and variation. Separate persistence, keepalive, and time-wait concerns prevent deadlock or stale state in specific cases.\item \textbf{6.5.10 Congestion control} Slow start grows the congestion window rapidly from a small value; congestion avoidance grows more cautiously; loss triggers a reduction. Duplicate acknowledgments can support fast retransmission and recovery.\item \textbf{6.5.11 Future of TCP} Evolving demands motivate larger windows, timestamps, selective acknowledgments, and improved congestion algorithms while preserving interoperability.\end{itemize}\end{minipage} \\ \hline
\cellcolor{CCue}\textbf{6.6 Performance issues} & \begin{minipage}[t]{\linewidth}\raggedright \begin{itemize}\item \textbf{6.6.1 Performance problems} Bottlenecks may arise in links, routers, host operating systems, memory movement, protocol processing, or application design. Faster links alone do not guarantee faster transfers.\item \textbf{6.6.2 Measurement} Measure the relevant workload end to end, account for cache and startup effects, vary one factor carefully, and distinguish throughput, latency, loss, and CPU cost.\item \textbf{6.6.3 Host design} Efficient hosts reduce copying, context switches, interrupts, and repeated traversal of protocol data; integrated processing and adequate buffering sustain fast links.\item \textbf{6.6.4 Fast segment processing} Optimize the common established-data path and handle exceptional cases separately, while preserving correctness.\item \textbf{6.6.5 Header compression} On low-bandwidth links, predictable header fields can be represented by small changes relative to prior packets, trading state and error sensitivity for lower overhead.\item \textbf{6.6.6 Long fat networks} A path with large bandwidth-delay product requires large windows and enough sequence space; timestamps and scaling prevent throughput from being limited by small legacy fields.\end{itemize}\end{minipage} \\ \hline
\cellcolor{CCue}\textbf{6.7 Delay-tolerant networking} & \begin{minipage}[t]{\linewidth}\raggedright \begin{itemize}\item \textbf{6.7.1 Architecture} When an end-to-end path may not exist, nodes store data persistently and forward it opportunistically across regions with long delay or disruption.\item \textbf{6.7.2 Bundle Protocol} Bundles provide an overlay data unit that can be held for long periods, transferred between regional transports, and managed with endpoint and custody concepts suited to disruption.\end{itemize}\end{minipage} \\ \hline
\cellcolor{CCue}\textbf{6.8 Summary} & \begin{minipage}[t]{\linewidth}\raggedright Transport protocols create process-to-process service above packet delivery. Endpoint naming, freshness, reliable release, windows, timers, and congestion response must agree; performance depends on both path properties and host implementation.\end{minipage} \\ \hline
\end{longtable}

\begin{legacybox}The sockets interface, classic TCP algorithms, header-compression discussion, and delay-tolerant overlay are preserved in the source's historical framing. Classic algorithm descriptions are foundations, not a complete catalog of present TCP implementations.\end{legacybox}

\notesection{Terminology and notation}

\begin{longtable}{@{}L{0.245\textwidth}|L{0.695\textwidth}@{}}
\rowcolor{CNavy}\color{white}\bfseries Cue / recall prompt &
\color{white}\bfseries Notes \\
\endfirsthead
\rowcolor{CNavy}\color{white}\bfseries Cue / recall prompt &
\color{white}\bfseries Notes (continued) \\
\endhead
\hline
\endfoot
\cellcolor{CCue}\textbf{TSAP / port} & \begin{minipage}[t]{\linewidth}\raggedright Transport endpoint identifier used to direct data to a process.\end{minipage} \\ \hline
\cellcolor{CCue}\textbf{Socket} & \begin{minipage}[t]{\linewidth}\raggedright Local software endpoint and programming abstraction for network communication.\end{minipage} \\ \hline
\cellcolor{CCue}\textbf{Three-way handshake} & \begin{minipage}[t]{\linewidth}\raggedright Exchange that confirms fresh bidirectional state and initial sequence values.\end{minipage} \\ \hline
\cellcolor{CCue}\textbf{Receive window} & \begin{minipage}[t]{\linewidth}\raggedright Receiver-advertised amount of data that may be outstanding without overrun.\end{minipage} \\ \hline
\cellcolor{CCue}\textbf{Congestion window} & \begin{minipage}[t]{\linewidth}\raggedright Sender-maintained limit reflecting estimated network capacity.\end{minipage} \\ \hline
\cellcolor{CCue}\textbf{RTT} & \begin{minipage}[t]{\linewidth}\raggedright Round-trip time from sending data until the corresponding feedback returns.\end{minipage} \\ \hline
\cellcolor{CCue}\textbf{RTO} & \begin{minipage}[t]{\linewidth}\raggedright Retransmission timeout derived from round-trip behavior.\end{minipage} \\ \hline
\cellcolor{CCue}\textbf{Byte stream} & \begin{minipage}[t]{\linewidth}\raggedright Ordered sequence of bytes without preserved application record boundaries.\end{minipage} \\ \hline
\cellcolor{CCue}\textbf{RPC} & \begin{minipage}[t]{\linewidth}\raggedright Remote procedure call represented through marshaled request and reply messages.\end{minipage} \\ \hline
\cellcolor{CCue}\textbf{Bandwidth-delay product} & \begin{minipage}[t]{\linewidth}\raggedright Data required in flight to fill a path: bandwidth multiplied by round-trip delay.\end{minipage} \\ \hline
\end{longtable}

\notesection{Comparison matrix}

\begin{longtable}{@{}L{0.313\textwidth}L{0.313\textwidth}L{0.313\textwidth}@{}}
\toprule
\textbf{Protocol / mechanism} & \textbf{Provides} & \textbf{Does not itself provide} \\ \midrule
UDP & Ports, datagrams, length, checksum & Connection, ordering, retransmission, or congestion control \\
RTP over a transport & Media sequence, timing, type, and source information & Resource reservation or guaranteed delivery \\
TCP & Reliable ordered full-duplex byte stream with flow and congestion control & Application message boundaries \\
Receive window & Protection against receiver-buffer overrun & Protection of network queues \\
Congestion window & Limit based on perceived network conditions & A promise of available bandwidth \\
\bottomrule
\end{longtable}
\notesection{Chapter synthesis}

\begin{overviewbox}\textbf{Cause-and-effect chain.} Applications need process delivery; ports name processes; unreliable packet service creates loss, duplication, and reordering; handshakes establish fresh state; acknowledgments, timers, and windows provide reliability and flow regulation; shared network capacity requires congestion control; large products of rate and delay require larger windows and careful host processing.\par
\textbf{Common confusions.} TCP is a byte stream rather than a message service; its receive window and congestion window solve different problems; UDP's checksum does not make delivery reliable; and an RPC timeout does not reveal whether the remote operation executed.\par
\textbf{Derived insight.} End-to-end correctness sometimes extends above transport: only the application can safely reconcile externally visible operations after ambiguous failures.\end{overviewbox}

\notesection{Retrieval practice}

\begin{enumerate}
  \item Trace the socket calls made by a connection-oriented server and client.
  \item Why are port numbers needed in addition to network addresses?
  \item How does a three-way handshake reject delayed duplicate connection requests?
  \item Why can perfect connection release not be guaranteed over an unreliable channel?
  \item Distinguish upward from downward multiplexing.
  \item Why can a timeout not give RPC exactly-once knowledge by itself?
  \item What information does RTP add for real-time media, and what does it not guarantee?
  \item Trace TCP establishment and release through their key states.
  \item Why does TCP acknowledge bytes rather than application messages?
  \item How do receive and congestion windows jointly bound a sender?
  \item Why must retransmission timeout account for RTT variation?
  \item Describe slow start and congestion avoidance after acknowledgment and loss events.
  \item Calculate the window needed to fill a path from its rate and RTT.
  \item Why is persistent storage central to delay-tolerant networking?
\end{enumerate}
\begin{CornellSummaryBox}{Rapid-review Cornell summary}Transport names processes and builds end-to-end service above packets. UDP preserves datagrams with minimal control; RTP adds media timing information; TCP adds a reliable ordered byte stream. Handshakes establish fresh state, windows handle receiver and network limits, timers recover loss, and performance depends on the bandwidth-delay product and host processing.\end{CornellSummaryBox}
\end{document}
